\documentclass[11pt]{article}

\usepackage{fullpage}

\begin{document}

\title{ARM Interim Checkpoint}
\author{Tangina Kader, Amy Showell, Shriya Sistla, Yao Tong}

\maketitle

\section{Group Organisation}

The group is organized such that all members receive a similar amount of work, and 
each member takes ownership of their own emulator tasks. To assign these tasks, 
we first gathered in-person to read the specification and create a Google Docs checklist of emulator components as 
extracted from the spec. Then, each member of the group volunteered to complete a few tasks they found interesting. Since 
it is hard to estimate exactly how long each task will take, some leftover tasks were kept open 
to be claimed as group members completed their initial assignments. 

We coordinate our work by communicating progress and requests through a WhatsApp 
channel. For example, when a group member finishes one of their tasks, they typically announce 
the completed task and asks others to review the merge request. After the request 
is reviewed and adjustments are made, the code is merged and a reminder to pull is sent to the 
channel. Lastly, the group member claims an unfinished task from the checklist by writing 
their name next to an un-claimed task, allowing us to easily track progress. Although each member 
works independently, sometimes there are sections of code written by others that are relevant 
yet need to be updated, patched, or pushed. These are also communicated and documented clearly 
in the WhatsApp channel. 

A good combination of remote and in-person work helps the group function exceptionally well. When working 
from home, it is easier to concentrate on and complete the task at hand. When working in-person during lab 
sessions, it is easier to patch bugs and communicate concerns. As a result, all the 
emulator tasks were completed ahead of the schedule we had made on day one. 

Although communication is our strength, one aspect the group can improve on is communicating potential 
utility functions -- like bit masks and memory reads -- so they can be generalized immediately instead 
of at the end, saving time and energy in the long-term. In the planning stages, the group should improve on
identifying closely-related tasks. Through identifying these tasks and assigning them to the same person, 
we can save time by having one group member complete the entire component in one branch instead of completing
half a component, reviewing, merging, and having a different member continue where the previous left off.  

\section{Implementation Strategies}

Our emulator is split into four main sections: the "hardware" definitions, the initialization and fetch cycle, the execution, 
and the global utility functions. First, the \texttt{PSTATE} register is represented by a struct consisting of four Booleans, each 
representing a \texttt{PSTATE} flag. The state of the ARM is represented by a struct consisting of an 8-bit int array representing memory,
a 64-bit int array representing the 31 general-purpose registers, a 64-bit int representing the PC, and the \texttt{PSTATE} struct 
representing the Processor State Register. These components are defined in their own armv8.h file, and can be updated using a 
subset of the global utility functions, "modify-regs". Note that the zero register is not included, since it represents the 
value zero and so can be handled independently. 

Next, the ARM memory is initialized using \texttt{memset()}, setting 2MB of memory to zero before instructions and data
are loaded. This is freed at the end of execution. The remaining components are initialized to their defaults as stated
in the specification file. The fetch cycle repeatedly calls a "decode" function that reads the instruction corresponding 
to the address in the PC, then passes the instruction to the proper function for execution. This decode function then 
returns a status code indicating successful execution, failed execution, or halt condition reached. 

The execution of instructions is divided into three modules: Data Processing, Data Transfer, and Branch. These are accessed by 
decode depending on bits 28-25 of each instruction, and all return status codes representing success or failure of execution.

Lastly, the global utility functions are composed of functions to read and write registers, update the PSTATE register, and
perform sign extensions. These are frequently used across all execution modules, and serve to reduce clutter in the execution 
modules' code. 

Some implementation tasks we will likely find challenging are identifying utility functions early on and maintaining code organization. 
As mentioned previously, some of our utility functions were refactored after the emulator functionality was complete, which is 
time-consuming and inefficient. To mitigate this challenge, group members should use the group chat to ask 
if someone already has an implementation for a problem that seems very common, such as extracting bits from an instruction or 
sign-extending an offset. This will speed development by reducing redundancy, shorten source files, and reduce errors. 

\end{document}
